\chapter[11]{Closed Manifolds with Positive Curvature}
\label{ch:chow-ii-closed-positive-curvature}

Throughout this chapter, \(E=\mathbb R^n\) is Euclidean space,
\(N=n(n-1)/2\), and

\[
  \mathcal A_n=\{R\in S^2(\Lambda^2E):R\text{ satisfies the first Bianchi identity}\}
\]

is the space of algebraic curvature operators.  We write
\(Q(R)=R^2+R^\#\), where \(\#\) is Hamilton's sharp product.

\section{Curvature operators and their algebra}

\begin{definition}[two-positive curvature operator]
  \label{def:chow-ii-two-positive-curvature-operator}
  \source{chow-et-al-ricci-flow-part-ii:page-0094}
  \dcref{ch11:11.1}
  \group{Positive Curvature Operators}
  A Riemannian manifold has \emph{two-positive curvature operator} if every
  sum of two distinct eigenvalues of its Riemann curvature operator is
  positive.  Replacing ``positive'' by ``nonnegative'' defines a
  \emph{two-nonnegative curvature operator}.
\end{definition}

\begin{definition}[sharp product and curvature quadratic]
  \label{def:chow-ii-sharp-product}
  \source{chow-et-al-ricci-flow-part-ii:page-0099,chow-et-al-ricci-flow-part-ii:page-0100}
  \group{Curvature Operator Algebra}
  Let \(\{\omega_\alpha\}\) be an orthonormal basis of \(\mathfrak{so}(n)\),
  with structure constants
  \([\omega_\alpha,\omega_\beta]=c_{\alpha\beta}^{\gamma}\omega_\gamma\).
  For \(A,B\in S^2(\mathfrak{so}(n))\), define
  \[
    (A\#B)_{\alpha\beta}
      =\frac12c_{\alpha}^{\gamma\theta}c_{\beta}^{\varepsilon\zeta}
        A_{\gamma\varepsilon}B_{\theta\zeta},
    \qquad Q(A)=A^2+A^\#,
  \]
  where \(A^\#=A\#A\).  Its polarization is
  \(Q(A,B)=\tfrac12(AB+BA)+A\#B\).
\end{definition}

\begin{lemma}[intrinsic formula for the sharp product]
  \label{lem:chow-ii-sharp-product-intrinsic-formula}
  \source{chow-et-al-ricci-flow-part-ii:page-0100}
  \uses{def:chow-ii-sharp-product}
  \dcref{ch11:11.3}
  \group{Curvature Operator Algebra}
  For every orthonormal basis \(\{\omega_\alpha\}\) of \(\Lambda^2E\),
  \[
  \langle(A\#B)\phi,\psi\rangle
  =\frac12\sum_{\alpha,\beta}
    \langle[A\omega_\alpha,B\omega_\beta],\phi\rangle
    \langle[\omega_\alpha,\omega_\beta],\psi\rangle .
  \]
  Consequently \(\#\) is symmetric, bilinear, and \(O(n)\)-equivariant.
\end{lemma}
\begin{proof}
\sketch
  Expand both brackets in structure constants.  The coefficient of
  \(\phi^\eta\psi^\theta\) is
  \(\tfrac12c_\eta^{\gamma\delta}c_\theta^{\alpha\beta}
  A_{\gamma\alpha}B_{\delta\beta}\), which is the defining coefficient of
  \(A\#B\).  Bilinearity is immediate.  Total antisymmetry of the structure
  constants gives symmetry in \(A,B\), and invariance of the bracket and
  inner product under \(O(n)\) gives equivariance.
\end{proof}

\section{The cone of two-nonnegative operators}

\begin{definition}[two-nonnegative algebraic curvature operator]
  \label{def:chow-ii-two-nonnegative-algebraic-curvature}
  \source{chow-et-al-ricci-flow-part-ii:page-0133}
  \dcref{ch11:11.37}
  \group{Positive Curvature Operators}
  An operator \(R\in\mathcal A_n\) is \emph{two-nonnegative} if the sum of
  every two distinct eigenvalues is nonnegative, and \emph{two-positive} if
  all these sums are positive.  Write \(\mathcal C\) for the closed convex
  cone of two-nonnegative operators.
\end{definition}

\begin{definition}[geometrically nonnegative curvature operator]
  \label{def:chow-ii-geometrically-nonnegative-curvature}
  \source{chow-et-al-ricci-flow-part-ii:page-0133,chow-et-al-ricci-flow-part-ii:page-0134}
  \dcref{ch11:11.38}
  \group{Positive Curvature Operators}
  An algebraic curvature operator is \emph{geometrically nonnegative} if it
  is a finite sum \(\sum_\alpha c_\alpha\omega_\alpha\otimes\omega_\alpha\),
  where \(c_\alpha>0\) and every
  \(\omega_\alpha\otimes\omega_\alpha\) satisfies the first Bianchi identity.
\end{definition}

\begin{lemma}[elementary consequences of two-nonnegativity]
  \label{lem:chow-ii-two-nonnegative-consequences}
  \source{chow-et-al-ricci-flow-part-ii:page-0134,chow-et-al-ricci-flow-part-ii:page-0135}
  \uses{def:chow-ii-two-nonnegative-algebraic-curvature, def:chow-ii-geometrically-nonnegative-curvature, def:chow-ii-algebraic-ricci-contraction}
  \dcref{ch11:11.39}
  \group{Positive Curvature Operators}
  Let \(\mu_1\le\cdots\le\mu_N\) be the eigenvalues of a two-nonnegative
  \(R\).  Then
  \[
    \mu_j\ge0\ (j\ge2),\qquad
    \mu_1\ge-\mu_2\ge-\frac{\operatorname{tr}(R)}{N-2}.
  \]
  Moreover \(\operatorname{Scal}(R)\ge0\), with equality only for \(R=0\),
  and \(\operatorname{Rc}(R)\ge0\).  The cone \(\mathcal C\) contains every
  geometrically nonnegative operator.
\end{lemma}
\begin{proof}
\sketch
  Two-nonnegativity permits at most one negative eigenvalue and gives
  \(-\mu_1\le\mu_2\).  Also
  \(\operatorname{tr}(R)\ge\mu_1+\mu_2+(N-2)\mu_2\ge(N-2)\mu_2\), proving
  the eigenvalue bounds.  Summing all inequalities \(\mu_i+\mu_j\ge0\)
  proves nonnegative scalar contraction, and equality forces all eigenvalues
  to vanish.  For a unit vector \(v\), the Ricci curvature is the trace of
  \(R\) on the \((n-1)\)-plane \(v\wedge v^\perp\); the min--max principle
  bounds it below by \(\mu_1+\cdots+\mu_{n-1}\ge0\).  Finally, each summand in
  a geometrically nonnegative operator is positive semidefinite, so their sum
  belongs to \(\mathcal C\).
\end{proof}

\begin{lemma}[geometrically nonnegative defect terms]
  \label{lem:chow-ii-geometric-nonnegativity-defect}
  \source{chow-et-al-ricci-flow-part-ii:page-0135,chow-et-al-ricci-flow-part-ii:page-0136}
  \uses{def:chow-ii-geometrically-nonnegative-curvature, thm:chow-ii-conjugation-defect-formula}
  \dcref{ch11:11.40}
  \group{Invariant Curvature Cones}
  If \(A\ge0\), then \(A\wedge A\), \(A_0^2\wedge\operatorname{id}\), and
  \(I\) are geometrically nonnegative.  Consequently, if
  \[
  0\le b\le\frac{\sqrt{2n(n-2)+4}-2}{n(n-2)},
  \qquad 2a=2b+(n-2)b^2,
  \]
  and \(\operatorname{Rc}(R)\ge0\), then \(D_{a,b}(R)\) is geometrically
  nonnegative.
\end{lemma}
\begin{proof}
\sketch
  Diagonalize \(A\), with eigenvalues \(\rho_i\ge0\), and diagonalize
  \(A_0\), with eigenvalues \(\lambda_i\).  On \(e_i\wedge e_j\), the three
  operators in the first assertion have eigenvalues
  \(\rho_i\rho_j\), \((\lambda_i^2+\lambda_j^2)/2\), and \(1\), respectively;
  each is therefore a nonnegative sum of decomposable rank-one operators.
  Under the displayed relation the first term in
  \ref{thm:chow-ii-conjugation-defect-formula} vanishes and the coefficient
  of \(I\) becomes
  \[
  \frac{\sigma b^2}{1+2(n-1)a}\bigl(2-4b-n(n-2)b^2\bigr)\ge0.
  \]
  The remaining terms have the forms just considered, with
  \(A=\operatorname{Rc}(R)\).
\end{proof}

\begin{proposition}[the curvature ODE preserves two-nonnegativity]
  \label{prop:chow-ii-ode-preserves-two-nonnegative}
  \source{chow-et-al-ricci-flow-part-ii:page-0136,chow-et-al-ricci-flow-part-ii:page-0137}
  \uses{def:chow-ii-two-nonnegative-algebraic-curvature, def:chow-ii-sharp-product}
  \dcref{ch11:11.42}
  \group{Invariant Curvature Cones}
  The cone \(\mathcal C\) is preserved by \(\dot R=Q(R)\).
\end{proposition}
\begin{proof}
\sketch
  Diagonalize a boundary point \(R\), with eigenvalues
  \(\mu_1\le\cdots\le\mu_N\) and eigenvectors \(\omega_\alpha\).  It suffices
  to prove \(Q(R)_{\alpha\alpha}+Q(R)_{\beta\beta}\ge0\) whenever
  \(\mu_\alpha+\mu_\beta=0\).  Unless \(R\ge0\), all such pairs are
  \((1,j)\), with \(-\mu_1=\mu_j>0\).  The structure-constant formula gives
  \[
  Q(R)_{11}+Q(R)_{jj}=\mu_1^2+\mu_j^2+
  2\sum_{\beta<\gamma}\bigl((c_1^{\beta\gamma})^2+
  (c_j^{\beta\gamma})^2\bigr)\mu_\beta\mu_\gamma.
  \]
  Pair every possibly negative term, which contains \(\mu_1\), with the
  term having the same squared structure constant and factor \(\mu_j\).
  Their sum is \((\mu_1+\mu_j)\mu_\gamma=0\); all unpaired terms are
  nonnegative.  Thus \(Q(R)\) belongs to the tangent cone at every boundary
  point, proving invariance.
\end{proof}

\begin{proposition}[Ricci flow preserves two-nonnegative curvature]
  \label{prop:chow-ii-ricci-flow-preserves-two-nonnegative}
  \source{chow-et-al-ricci-flow-part-ii:page-0137,chow-et-al-ricci-flow-part-ii:page-0138}
  \uses{prop:chow-ii-curvature-operator-evolution, prop:chow-ii-ode-preserves-two-nonnegative}
  \dcref{ch11:11.44}
  \group{Ricci Flow Maximum Principles}
  A complete bounded-curvature Ricci flow whose initial curvature operator is
  two-nonnegative has two-nonnegative curvature operator at every later time.
\end{proposition}
\begin{proof}
\sketch
  Uhlenbeck's bundle identification converts the curvature equation into
  \((\partial_t-\Delta)\operatorname{Rm}=Q(\operatorname{Rm})\).  The fiberwise
  cone corresponding to \(\mathcal C\) is closed, convex, invariant under
  parallel translation, and, by
  \ref{prop:chow-ii-ode-preserves-two-nonnegative}, invariant under the
  reaction ODE.  The tensor maximum principle therefore keeps
  \(\operatorname{Rm}\) in this cone for all later times.
\end{proof}

\begin{proposition}[strong alternative for two-nonnegative curvature]
  \label{prop:chow-ii-strong-two-nonnegative-alternative}
  \source{chow-et-al-ricci-flow-part-ii:page-0138,chow-et-al-ricci-flow-part-ii:page-0139,chow-et-al-ricci-flow-part-ii:page-0140}
  \uses{prop:chow-ii-ricci-flow-preserves-two-nonnegative, prop:chow-ii-ode-preserves-two-nonnegative}
  \dcref{ch11:11.45}
  \group{Ricci Flow Maximum Principles}
  For a complete bounded-curvature Ricci flow initially two-nonnegative, at
  every positive time the curvature operator is either nonnegative or
  two-positive.  If it is two-positive at one space-time point, then it is
  two-positive everywhere at every later time.
\end{proposition}
\begin{proof}
\sketch
  Let \(u=\mu_1+\mu_2\ge0\).  On a relatively compact domain, compare the
  curvature operator with
  \(\operatorname{Rm}+(\varepsilon e^{At}-f)I\), where
  \(f_t=\Delta f-Af\), \(f\) vanishes on the boundary, and its initial value
  is nonnegative and positive at the given point.  For \(A\) large the
  perturbed operator is a strict supersolution; the tensor maximum principle
  and then \(\varepsilon\downarrow0\) give \(u>0\) throughout the domain at
  later times.  If two-positivity never occurs, then \(u\equiv0\).  At a
  minimum, the tangent-cone computation in
  \ref{prop:chow-ii-ode-preserves-two-nonnegative} gives
  \(0\ge\mu_1^2+\mu_2^2\), hence \(\mu_1=\mu_2=0\) and the operator is
  nonnegative.  Exhausting the complete manifold proves both assertions.
\end{proof}

\begin{proposition}[preserved stretched two-nonnegative cones]
  \label{prop:chow-ii-preserved-stretched-two-nonnegative-cones}
  \source{chow-et-al-ricci-flow-part-ii:page-0140,chow-et-al-ricci-flow-part-ii:page-0141,chow-et-al-ricci-flow-part-ii:page-0142}
  \uses{lem:chow-ii-linear-image-ode-cone-criterion, lem:chow-ii-geometric-nonnegativity-defect, prop:chow-ii-ode-preserves-two-nonnegative}
  \dcref{ch11:11.47}
  \group{Invariant Curvature Cones}
  Let \(n\ge4\),
  \[
  \bar b=\frac{\sqrt{2n(n-2)+4}-2}{n(n-2)},\qquad
  0\le b\le\bar b,\qquad 2a=2b+(n-2)b^2.
  \]
  Then \(\ell_{a,b}(\mathcal C)\) is preserved by \(\dot R=Q(R)\).  For
  \(b>0\), \(Q\) is transverse to its boundary away from the origin, and
  \(\ell_{a,\bar b}(\mathcal C)\setminus\{0\}\) consists of positive
  curvature operators.
\end{proposition}
\begin{proof}
\sketch
  Pull the vector field back to \(\mathcal C\).  It is
  \(X_{a,b}=Q+D_{a,b}\).  The first summand is tangent by
  \ref{prop:chow-ii-ode-preserves-two-nonnegative}; the second is
  geometrically nonnegative by
  \ref{lem:chow-ii-geometric-nonnegativity-defect}, hence belongs to every
  tangent cone of \(\mathcal C\).  This proves invariance by
  \ref{lem:chow-ii-linear-image-ode-cone-criterion}.  If \(0<b<\bar b\),
  the identity coefficient in \(D_{a,b}\) is positive unless
  \(\operatorname{Rc}_0=0\); in that case the positive scalar part of
  \(2a\operatorname{Rc}\wedge\operatorname{Rc}\) is positive.  This gives
  transversality.  At \(b=\bar b\), use
  \(\operatorname{Rc}(R)\ge(n-3)|\mu_1(R)|\operatorname{id}\) from
  \ref{lem:chow-ii-two-nonnegative-consequences} in the formula for
  \(\ell_{a,b}(R)\); the identity
  \(2(1-2\bar b)=n(n-2)\bar b^2\) makes its smallest eigenvalue positive.
  This proves the endpoint claim and, by continuity, endpoint transversality.
\end{proof}

\section{Pinching families and convergence}

\begin{definition}[pinching family]
  \label{def:chow-ii-pinching-family}
  \source{chow-et-al-ricci-flow-part-ii:page-0143}
  \dcref{ch11:11.50}
  \group{Ricci Flow Pinching}
  A continuous family \(C(s)\subset\mathcal A_n\), \(0\le s<1\), of
  full-dimensional closed convex cones is a \emph{pinching family} when each
  cone is \(O(n)\)-invariant, every nonzero member has positive scalar
  curvature, \(Q(R)\) lies in the interior of the tangent cone at every
  nonzero boundary point for \(0<s<1\), and \(C(s)\) converges in pointed
  Hausdorff topology to \(\mathbb R_+I\) as \(s\to1\).
\end{definition}

\begin{theorem}[pinching family from nonnegative curvature]
  \label{thm:chow-ii-pinching-family-nonnegative-curvature}
  \source{chow-et-al-ricci-flow-part-ii:page-0143,chow-et-al-ricci-flow-part-ii:page-0144,chow-et-al-ricci-flow-part-ii:page-0145,chow-et-al-ricci-flow-part-ii:page-0146,chow-et-al-ricci-flow-part-ii:page-0147,chow-et-al-ricci-flow-part-ii:page-0148,chow-et-al-ricci-flow-part-ii:page-0149,chow-et-al-ricci-flow-part-ii:page-0150,chow-et-al-ricci-flow-part-ii:page-0151,chow-et-al-ricci-flow-part-ii:page-0152,chow-et-al-ricci-flow-part-ii:page-0153}
  \uses{def:chow-ii-pinching-family, cor:chow-ii-conjugation-defect-eigenvalues, lem:chow-ii-linear-image-ode-cone-criterion}
  \dcref{ch11:11.51}
  \group{Ricci Flow Pinching}
  There is a pinching family whose initial cone is the cone of nonnegative
  algebraic curvature operators.
\end{theorem}
\begin{proof}
\sketch
  For \(0\le b\le1/2\), put
  \[
  p_1(b)=\frac{(n-2)b^2}{1+(n-2)b^2},\quad
  a_1(b)=\frac{(n-2)b^2+2b}{2+2(n-2)b^2},
  \]
  and let \(C_b\) consist of operators \(R\ge0\) satisfying
  \(\operatorname{Rc}(R)\ge p_1(b)\operatorname{Scal}(R)\operatorname{id}/n\).
  The eigenvalue formula in
  \ref{cor:chow-ii-conjugation-defect-eigenvalues}, together with the boundary
  equality \(\lambda_i=-(1-p_1)\bar\lambda\), shows that
  \(X_{a_1(b),b}\) is strictly inside both tangent half-spaces defining
  \(C_b\) away from zero.  Indeed its curvature eigenvalues are sums of
  \(2a_1(\lambda_i+(1-p_1)\bar\lambda)
  (\lambda_j+(1-p_1)\bar\lambda)/(1-p_1)\) and nonnegative square terms;
  its Ricci boundary excess is at least
  \(2nb^2\sigma/(nb+1)>0\).  Hence
  \(\ell_{a_1(b),b}(C_b)\) is preserved and transverse.

  Next, for \(s\ge0\), put
  \(a_2(s)=(1+s)/2\), \(p_2(s)=1-4/(n+2+4s)\), and define \(C'_s\) by the
  analogous nonnegative-curvature and Ricci-pinching inequalities.  The same
  eigenvalue formula, using
  \(\sigma\le(n-1)(1-p_2)^2\bar\lambda^2\), proves strict positivity of
  \(X_{a_2(s),1/2}\) and strict improvement of the Ricci boundary inequality.
  Thus \(\ell_{a_2(s),1/2}(C'_s)\) is preserved and transverse.  The first
  family at \(b=1/2\) equals the second at \(s=0\).  Concatenate them after
  reparametrizing \(s\in[0,\infty)\) to \([1/2,1)\).  Since
  \(a_2(s)^{-1}\ell_{a_2(s),1/2}(R)\to2(n-1)R_I\) and the Ricci pinching tends
  to equality, the cones converge to \(\mathbb R_+I\).  All conditions in
  \ref{def:chow-ii-pinching-family} follow.
\end{proof}

\begin{corollary}[pinching family from two-nonnegative curvature]
  \label{cor:chow-ii-pinching-family-two-nonnegative}
  \source{chow-et-al-ricci-flow-part-ii:page-0143,chow-et-al-ricci-flow-part-ii:page-0144}
  \uses{prop:chow-ii-preserved-stretched-two-nonnegative-cones, thm:chow-ii-pinching-family-nonnegative-curvature}
  \dcref{ch11:11.52}
  \group{Ricci Flow Pinching}
  There is a pinching family whose initial cone is \(\mathcal C\), the cone of
  two-nonnegative algebraic curvature operators.
\end{corollary}
\begin{proof}
\sketch
  Let \(\bar b\) be as in
  \ref{prop:chow-ii-preserved-stretched-two-nonnegative-cones} and intersect
  every cone of the family in
  \ref{thm:chow-ii-pinching-family-nonnegative-curvature} with
  \(\ell_{a(\bar b),\bar b}(\mathcal C)\).  This remains a pinching family and
  begins at the latter cone.  Concatenate it with the preserved transverse
  family \(\ell_{a(b),b}(\mathcal C)\), \(0\le b\le\bar b\).
\end{proof}

\begin{definition}[pinching set]
  \label{def:chow-ii-pinching-set}
  \source{chow-et-al-ricci-flow-part-ii:page-0153,chow-et-al-ricci-flow-part-ii:page-0154}
  \dcref{ch11:11.59}
  \group{Ricci Flow Pinching}
  A subset \(Z\subset\mathcal A_n\) is a \emph{pinching set} if it is closed,
  convex, \(O(n)\)-invariant, preserved by \(\dot R=Q(R)\), and there are
  \(\delta>0\), \(K<\infty\) such that
  \[
    \left|R-\frac{\operatorname{tr}(R)}N I\right|
       \le K|R|^{1-\delta}\qquad(R\in Z).
  \]
\end{definition}

\begin{theorem}[generalized pinching set]
  \label{thm:chow-ii-generalized-pinching-set}
  \source{chow-et-al-ricci-flow-part-ii:page-0155,chow-et-al-ricci-flow-part-ii:page-0156,chow-et-al-ricci-flow-part-ii:page-0157,chow-et-al-ricci-flow-part-ii:page-0158,chow-et-al-ricci-flow-part-ii:page-0159,chow-et-al-ricci-flow-part-ii:page-0160,chow-et-al-ricci-flow-part-ii:page-0161}
  \uses{def:chow-ii-pinching-family}
  \dcref{ch11:11.64}
  \group{Ricci Flow Pinching}
  Let \(C(s)\) be a pinching family.  For every \(\varepsilon\in(0,1)\) and
  \(h_0>0\), there is a closed convex \(O(n)\)-invariant set
  \(F\subset\mathcal A_n\), preserved by \(\dot R=Q(R)\), such that
  \[
  \overline{C(\varepsilon)\cap\{R:\operatorname{tr}R\le h_0\}}
  \subset F\subset C(\varepsilon)
  \]
  and \(\overline{F\setminus C(s)}\) is compact for every
  \(s\in[\varepsilon,1)\).
\end{theorem}
\begin{proof}
\sketch
  Intersect all closed convex invariant \(O(n)\)-sets satisfying the two
  inclusions; call the minimal intersection \(F\).  It is preserved because
  every set in the intersection is preserved.  Suppose
  \(F\setminus C(s)\) first becomes unbounded at \(s_0\).  For small
  \(\delta>0\), shrink \(C(s_0)\) on the trace-one slice by distance
  \(\delta\), obtaining a cone \(C_\delta(s_0)\).  Transversality and the
  quadratic homogeneity of \(Q\) imply that all sufficiently small such
  cones remain invariant.

  Translate one of them by \(-I\) and truncate it by
  \(\operatorname{tr}R\ge h\).  For \(h\) large this truncated shifted cone is
  invariant: on the trace face,
  \(\operatorname{tr}Q(R)=|\operatorname{Rc}(R)|^2>0\); on the conical face,
  the inward distance of \(Q(R+I)\) is quadratic in \(|R|\), whereas
  \(|Q(R+I)-Q(R)|\) is only linear.  Continuity of \(C(s)\) and the minimal
  choice of \(s_0\) then permit a large dilation of this truncated cone to
  replace the unbounded part of \(F\).  The union with the bounded trace
  sublevel is again closed, convex, invariant, and contains the required
  initial slice, hence equals \(F\) by minimality.  Its complement from
  \(C(s)\) is bounded for all \(s\) near \(s_0\), contradicting the choice of
  \(s_0\).  Thus all the asserted complements are compact.
\end{proof}

\begin{theorem}[Boehm--Wilking convergence]
  \label{thm:chow-ii-boehm-wilking-convergence}
  \source{chow-et-al-ricci-flow-part-ii:page-0095,chow-et-al-ricci-flow-part-ii:page-0156,chow-et-al-ricci-flow-part-ii:page-0157}
  \uses{def:chow-ii-two-positive-curvature-operator}
  \dcref{ch11:11.2}
  \group{Positive Curvature Ricci Flow}
  Let \((M^n,g_0)\) be closed with two-positive curvature operator.  The
  volume-normalized Ricci flow starting at \(g_0\) exists for all positive
  time and converges exponentially in every \(C^k\)-norm to a metric of
  constant positive sectional curvature.
\end{theorem}
\begin{proof}
\sketch
  \uses{cor:chow-ii-pinching-family-two-nonnegative,thm:chow-ii-generalized-pinching-set,prop:chow-ii-ricci-flow-preserves-two-nonnegative}
  Compactness and strict two-positivity place the initial curvature operators
  in a bounded slice of \(C(\varepsilon)\) for some member of the pinching
  family from \ref{cor:chow-ii-pinching-family-two-nonnegative}.  Apply
  \ref{thm:chow-ii-generalized-pinching-set} and the tensor maximum principle
  to keep the curvature in the corresponding generalized pinching set \(F\).
  Positive scalar curvature forces a finite singular time for the unnormalized
  flow.  Choose maximum-curvature points and parabolically rescale.  No local
  collapsing and Hamilton compactness give a complete nonflat ancient limit.
  Since \(F\setminus C(s)\) is bounded for every \(s<1\), every nonzero
  curvature operator of the limit belongs to all \(C(s)\), hence to
  \(\mathbb R_+I\).  Schur's lemma makes its positive sectional curvature
  spatially constant.  The standard stability theorem for a constant positive
  curvature Ricci-flow limit upgrades this sequential convergence to
  exponential convergence of the volume-normalized flow in every
  \(C^k\)-norm, and continuation gives existence for all normalized time.
\end{proof}

\begin{definition}[wedge product of symmetric endomorphisms]
  \label{def:chow-ii-symmetric-endomorphism-wedge}
  \source{chow-et-al-ricci-flow-part-ii:page-0101,chow-et-al-ricci-flow-part-ii:page-0102}
  \group{Curvature Operator Algebra}
  For \(A,B\in S^2(E)\), define an endomorphism of \(\Lambda^2E\) by
  \[
    (A\wedge B)(v\wedge w)
      =\frac12\bigl(Av\wedge Bw+Bv\wedge Aw\bigr).
  \]
\end{definition}

\begin{lemma}[elementary properties of the symmetric wedge]
  \label{lem:chow-ii-symmetric-wedge-properties}
  \source{chow-et-al-ricci-flow-part-ii:page-0101,chow-et-al-ricci-flow-part-ii:page-0102}
  \uses{def:chow-ii-symmetric-endomorphism-wedge}
  \dcref{ch11:11.5}
  \group{Curvature Operator Algebra}
  For \(A,B\in S^2(E)\), one has \(A\wedge B=B\wedge A\),
  \(\operatorname{id}\wedge\operatorname{id}=\operatorname{id}_{\Lambda^2E}\),
  and \(A\wedge B\) is self-adjoint.  Moreover,
  \[
  (A\wedge B)_{ijk\ell}
   =\frac12(A_{ik}B_{j\ell}+B_{ik}A_{j\ell}
                 -A_{i\ell}B_{jk}-B_{i\ell}A_{jk}).
  \]
\end{lemma}
\begin{proof}
\sketch
  The first two assertions follow directly from the definition.  Expanding
  the inner product of \((A\wedge B)(e_i\wedge e_j)\) with
  \(e_k\wedge e_\ell\) gives the displayed formula, which is unchanged when
  the pairs \((i,j)\) and \((k,\ell)\) are interchanged because \(A\) and
  \(B\) are self-adjoint.  Thus \(A\wedge B\) is self-adjoint.
\end{proof}

\begin{proposition}[irreducible decomposition of two-tensors]
  \label{prop:chow-ii-two-tensor-decomposition}
  \source{chow-et-al-ricci-flow-part-ii:page-0103,chow-et-al-ricci-flow-part-ii:page-0104}
  \dcref{ch11:11.7}
  \group{Orthogonal Representations}
  As an orthogonal \(O(n)\)-representation,
  \[
    E\otimes E=\Lambda^2E\oplus S^2_0E\oplus\mathbb R\operatorname{id}.
  \]
\end{proposition}
\begin{proof}
\sketch
  Every endomorphism is uniquely the sum of a skew-adjoint and a
  self-adjoint endomorphism, and every self-adjoint endomorphism is uniquely
  the sum of its trace-free part and a scalar multiple of the identity.
  These summands are pairwise orthogonal and invariant under conjugation by
  \(O(n)\).  Their irreducibility is the standard irreducibility of
  \(\Lambda^2E\) and \(S^2_0E\) as real orthogonal representations, with the
  scalar summand one-dimensional.
\end{proof}

\begin{definition}[Bianchi map and algebraic curvature operators]
  \label{def:chow-ii-bianchi-map}
  \source{chow-et-al-ricci-flow-part-ii:page-0108,chow-et-al-ricci-flow-part-ii:page-0109}
  \group{Algebraic Curvature Operators}
  For a covariant four-tensor \(R\), set
  \[
    b(R)(x,y,z,t)=\frac13\bigl(R(x,y,z,t)+R(y,z,x,t)+R(z,x,y,t)\bigr).
  \]
  An element of \(S^2(\Lambda^2E)\) is an \emph{algebraic curvature
  operator} when it lies in \(\ker b\); their space is denoted \(\mathcal A_n\).
\end{definition}

\begin{lemma}[orthogonal Bianchi decomposition]
  \label{lem:chow-ii-bianchi-orthogonal-decomposition}
  \source{chow-et-al-ricci-flow-part-ii:page-0108,chow-et-al-ricci-flow-part-ii:page-0109,chow-et-al-ricci-flow-part-ii:page-0110}
  \uses{def:chow-ii-bianchi-map}
  \dcref{ch11:11.11}
  \group{Algebraic Curvature Operators}
  The restriction of \(b\) to \(S^2(\Lambda^2E)\) is a self-adjoint
  projection onto \(\Lambda^4E\).  Consequently
  \[
    S^2(\Lambda^2E)=\mathcal A_n\oplus\Lambda^4E
  \]
  orthogonally.
\end{lemma}
\begin{proof}
\sketch
  The pair symmetries of a member of \(S^2(\Lambda^2E)\) show directly that
  \(b(R)\) is alternating in all four entries.  If \(R\) is already a
  four-form, cyclic permutation of its first three entries has positive sign,
  so \(b(R)=R\).  Finally, reindexing the finite orthonormal-basis sum defining
  the tensor inner product gives \(\langle b(R),S\rangle=\langle R,b(S)\rangle\).
  Thus \(b\) is a self-adjoint projection, and its kernel and image are the
  asserted orthogonal summands.
\end{proof}

\begin{lemma}[wedge constructions satisfy Bianchi]
  \label{lem:chow-ii-wedge-satisfies-bianchi}
  \source{chow-et-al-ricci-flow-part-ii:page-0110,chow-et-al-ricci-flow-part-ii:page-0111}
  \uses{def:chow-ii-bianchi-map, def:chow-ii-symmetric-endomorphism-wedge}
  \dcref{ch11:11.13}
  \group{Algebraic Curvature Operators}
  If \(A\in S^2(E)\), then both
  \(A\wedge\operatorname{id}\) and \(A\wedge A\) belong to \(\mathcal A_n\).
\end{lemma}
\begin{proof}
\sketch
  Substitute the component formula of
  \ref{lem:chow-ii-symmetric-wedge-properties} into the cyclic sum defining
  \(b\).  Every term cancels with the term obtained by the symmetry of \(A\),
  both for \(B=\operatorname{id}\) and for \(B=A\).
\end{proof}

\begin{definition}[Ricci and scalar contractions]
  \label{def:chow-ii-algebraic-ricci-contraction}
  \source{chow-et-al-ricci-flow-part-ii:page-0111}
  \group{Algebraic Curvature Operators}
  For \(R\in\mathcal A_n\), define
  \[
    \operatorname{Rc}(R)_{ij}=\sum_kR_{ikjk},
    \qquad \operatorname{Scal}(R)=\operatorname{tr}(\operatorname{Rc}(R)).
  \]
\end{definition}

\begin{lemma}[Ricci contraction of a wedge]
  \label{lem:chow-ii-ricci-contraction-wedge}
  \source{chow-et-al-ricci-flow-part-ii:page-0112}
  \uses{def:chow-ii-symmetric-endomorphism-wedge, def:chow-ii-algebraic-ricci-contraction}
  \dcref{ch11:11.14}
  \group{Algebraic Curvature Operators}
  If \(A,B\in S^2(E)\), then
  \[
  \operatorname{Rc}(A\wedge B)
   =\frac12\{A(\operatorname{tr}(B)\operatorname{id}-B)
             +B(\operatorname{tr}(A)\operatorname{id}-A)\}.
  \]
  In particular,
  \(\operatorname{Rc}(A\wedge A)=\operatorname{tr}(A)A-A^2\) and
  \(\operatorname{Rc}(A\wedge\operatorname{id})
    =\tfrac{n-2}{2}A+\tfrac{\operatorname{tr}(A)}2\operatorname{id}\).
\end{lemma}
\begin{proof}
\sketch
  Contract the first and third indices in the component formula of
  \ref{lem:chow-ii-symmetric-wedge-properties}.  The four resulting terms
  are respectively \(A\operatorname{tr}(B)\), \(-AB\),
  \(B\operatorname{tr}(A)\), and \(-BA\), with the common factor \(1/2\).
  The special cases follow by substitution.
\end{proof}

\begin{lemma}[wedge embedding and Ricci adjunction]
  \label{lem:chow-ii-wedge-embedding-ricci-adjoint}
  \source{chow-et-al-ricci-flow-part-ii:page-0112,chow-et-al-ricci-flow-part-ii:page-0113}
  \uses{lem:chow-ii-ricci-contraction-wedge, lem:chow-ii-wedge-satisfies-bianchi}
  \dcref{ch11:11.15}
  \group{Algebraic Curvature Operators}
  If \(n>2\), the map \(A\mapsto\operatorname{id}\wedge A\) from
  \(S^2(E)\) to \(\mathcal A_n\) is injective and is the adjoint of
  \(\operatorname{Rc}:\mathcal A_n\to S^2(E)\).
\end{lemma}
\begin{proof}
\sketch
  If \(\operatorname{id}\wedge A=0\), the scalar trace of
  \ref{lem:chow-ii-ricci-contraction-wedge} gives
  \((n-1)\operatorname{tr}(A)=0\), and the same formula then gives
  \((n-2)A/2=0\).  For adjunction, expand both inner products in an
  orthonormal basis and use the curvature symmetries:
  \(\langle\operatorname{id}\wedge A,R\rangle
  =\tfrac12R_{ikjk}A_{ij}=\langle A,\operatorname{Rc}(R)\rangle\), with the
  common normalization induced on the two tensor spaces.
\end{proof}

\begin{lemma}[scalar--Ricci--Weyl decomposition]
  \label{lem:chow-ii-algebraic-curvature-decomposition}
  \source{chow-et-al-ricci-flow-part-ii:page-0113,chow-et-al-ricci-flow-part-ii:page-0114}
  \uses{lem:chow-ii-wedge-embedding-ricci-adjoint, def:chow-ii-algebraic-ricci-contraction}
  \dcref{ch11:11.16}
  \group{Algebraic Curvature Operators}
  For \(n>2\), every \(R\in\mathcal A_n\) decomposes orthogonally and uniquely as
  \[
  R=\frac{\operatorname{Scal}(R)}{n(n-1)}I
    +\frac2{n-2}\operatorname{id}\wedge\operatorname{Rc}_0(R)+R_W,
  \]
  where \(I=\operatorname{id}\wedge\operatorname{id}\),
  \(\operatorname{Rc}_0=\operatorname{Rc}-\operatorname{Scal}\operatorname{id}/n\),
  and \(R_W\in\ker\operatorname{Rc}\).
\end{lemma}
\begin{proof}
\sketch
  By \ref{lem:chow-ii-wedge-embedding-ricci-adjoint}, the orthogonal
  complement of the image of \(A\mapsto\operatorname{id}\wedge A\) is
  \(\ker\operatorname{Rc}\).  Split \(S^2(E)\) into its scalar and trace-free
  summands and use \ref{lem:chow-ii-ricci-contraction-wedge} to invert the
  Ricci contraction on each: it multiplies the scalar identity by \(n-1\)
  and a trace-free tensor by \((n-2)/2\).  This gives the displayed formula
  and uniqueness.
\end{proof}

\begin{corollary}[decomposition of a trace-free square wedge]
  \label{cor:chow-ii-trace-free-square-wedge-decomposition}
  \source{chow-et-al-ricci-flow-part-ii:page-0114}
  \uses{lem:chow-ii-algebraic-curvature-decomposition, lem:chow-ii-ricci-contraction-wedge}
  \dcref{ch11:11.17}
  \group{Algebraic Curvature Operators}
  If \(A\in S^2(E)\) and \(\operatorname{tr}(A)=0\), then
  \[
  A\wedge A=-\frac{\operatorname{tr}(A^2)}{n(n-1)}I
    -\frac2{n-2}(A^2)_0\wedge\operatorname{id}+(A\wedge A)_W.
  \]
\end{corollary}
\begin{proof}
\sketch
  By \ref{lem:chow-ii-ricci-contraction-wedge},
  \(\operatorname{Rc}(A\wedge A)=-A^2\).  Insert this contraction and its
  trace into \ref{lem:chow-ii-algebraic-curvature-decomposition}.
\end{proof}

\begin{proposition}[curvature-operator evolution]
  \label{prop:chow-ii-curvature-operator-evolution}
  \source{chow-et-al-ricci-flow-part-ii:page-0105,chow-et-al-ricci-flow-part-ii:page-0106,chow-et-al-ricci-flow-part-ii:page-0107}
  \uses{def:chow-ii-sharp-product}
  \dcref{ch11:11.9}
  \group{Ricci Flow Evolution}
  Under Ricci flow, after identifying the evolving tangent bundles by
  Uhlenbeck's bundle isometries, the curvature operator satisfies
  \[
    (\partial_t-\Delta)\operatorname{Rm}
       =\operatorname{Rm}^2+\operatorname{Rm}^\#.
  \]
\end{proposition}
\begin{proof}
\sketch
  In an orthonormal frame the tensor evolution is
  \((\partial_t-\Delta)R_{abcd}
  =2(B_{abcd}-B_{badc}+B_{acbd}-B_{adbc})\), where
  \(B_{abcd}=-R_{apbq}R_{cpdq}\).  The first Bianchi identity gives
  \((\operatorname{Rm}^2)_{abcd}=2R_{abef}R_{cdef}\), which equals
  \(-4(B_{abcd}-B_{badc})\).  Expanding the intrinsic sharp formula from
  \ref{lem:chow-ii-sharp-product-intrinsic-formula} gives
  \((\operatorname{Rm}^\#)_{abcd}=4(B_{adbc}-B_{abcd})\).  Substitution,
  together with \(\operatorname{Rm}_{abcd}=-2R_{abcd}\), yields the equation.
\end{proof}

\begin{corollary}[Ricci trace of the curvature quadratic]
  \label{cor:chow-ii-curvature-quadratic-trace}
  \source{chow-et-al-ricci-flow-part-ii:page-0114,chow-et-al-ricci-flow-part-ii:page-0115}
  \uses{def:chow-ii-sharp-product, def:chow-ii-algebraic-ricci-contraction}
  \dcref{ch11:11.18}
  \group{Curvature Operator Algebra}
  For \(R\in\mathcal A_n\), one has \(Q(R)\in\mathcal A_n\) and
  \[
    \operatorname{Rc}(Q(R))_{ij}
      =\sum_{k,\ell}\operatorname{Rc}(R)_{k\ell}R_{ikj\ell},
    \qquad
    \operatorname{Scal}(Q(R))=|\operatorname{Rc}(R)|^2.
  \]
  In particular, \(\ker\operatorname{Rc}\) is invariant under \(Q\).
\end{corollary}
\begin{proof}
\sketch
  Expanding \(R^2+R^\#\) by the same component calculation as in
  \ref{prop:chow-ii-curvature-operator-evolution} shows first that its cyclic
  Bianchi sum vanishes.  Contracting its first and third indices, and then
  using the Bianchi identity
  \(R_{aebf}R_{cdef}=\tfrac12R_{abef}R_{cdef}\), gives the first displayed
  identity.  Taking its trace gives the second.  If \(\operatorname{Rc}(R)=0\),
  the first identity vanishes.
\end{proof}

\section{Equivariant changes of the curvature ODE}

\begin{definition}[the transformations \(\ell_{a,b}\)]
  \label{def:chow-ii-ell-transformations}
  \source{chow-et-al-ricci-flow-part-ii:page-0117}
  \uses{lem:chow-ii-algebraic-curvature-decomposition}
  \dcref{ch11:11.22}
  \group{Equivariant Curvature Transformations}
  For \(a,b\in\mathbb R\), define
  \[
   \ell_{a,b}(R)=R+2(n-1)aR_I+(n-2)bR_{\operatorname{Rc}_0}.
  \]
  Thus \(\ell_{a,b}\) multiplies the scalar, trace-free Ricci, and Weyl
  summands by \(1+2(n-1)a\), \(1+(n-2)b\), and \(1\), respectively.  It is
  invertible when the first two factors are nonzero.
\end{definition}

\begin{definition}[conjugated quadratic and defect]
  \label{def:chow-ii-conjugated-curvature-quadratic}
  \source{chow-et-al-ricci-flow-part-ii:page-0118}
  \uses{def:chow-ii-ell-transformations, def:chow-ii-sharp-product}
  \dcref{ch11:11.24}
  \group{Equivariant Curvature Transformations}
  For invertible \(\ell_{a,b}\), set
  \[
  X_{a,b}=\ell_{a,b}^{-1}\circ Q\circ\ell_{a,b},
  \qquad D_{a,b}=X_{a,b}-Q.
  \]
\end{definition}

\begin{lemma}[linear images and invariant ODE cones]
  \label{lem:chow-ii-linear-image-ode-cone-criterion}
  \source{chow-et-al-ricci-flow-part-ii:page-0118}
  \uses{def:chow-ii-conjugated-curvature-quadratic}
  \dcref{ch11:11.26}
  \group{Invariant Curvature Cones}
  Let \(C\subset\mathcal A_n\) be a closed convex \(O(n)\)-invariant cone.
  Then \(\ell_{a,b}(C)\) is preserved by \(\dot R=Q(R)\) if and only if
  \(C\) is preserved by \(\dot R=X_{a,b}(R)\).
\end{lemma}
\begin{proof}
\sketch
  A curve \(R(t)\) solves \(\dot R=Q(R)\) precisely when
  \(\ell_{a,b}^{-1}R(t)\) solves \(\dot S=X_{a,b}(S)\).  Since
  \(R(t)\in\ell_{a,b}(C)\) is equivalent to
  \(\ell_{a,b}^{-1}R(t)\in C\), invariance of either set is equivalent to
  invariance of the other.
\end{proof}

\begin{lemma}[interaction with the identity curvature operator]
  \label{lem:chow-ii-sharp-identity-formula}
  \source{chow-et-al-ricci-flow-part-ii:page-0126,chow-et-al-ricci-flow-part-ii:page-0127}
  \uses{def:chow-ii-sharp-product, def:chow-ii-algebraic-ricci-contraction}
  \dcref{ch11:11.33}
  \group{Curvature Operator Algebra}
  For \(R\in\mathcal A_n\),
  \[
    R+R\#I=\operatorname{Rc}(R)\wedge\operatorname{id}
      =(n-1)R_I+\frac{n-2}{2}R_{\operatorname{Rc}_0}.
  \]
\end{lemma}
\begin{proof}
\sketch
  In an orthonormal basis, the intrinsic sharp formula gives
  \[
  (R\#I)_{abcd}=\frac12(r_{ac}\delta_{bd}-r_{ad}\delta_{bc}
    +r_{bd}\delta_{ac}-r_{bc}\delta_{ad})-R_{abcd},
  \]
  where \(r=\operatorname{Rc}(R)\); the Bianchi identity is used to collect
  the four contraction terms.  The first equality follows from the component
  formula for \(r\wedge\operatorname{id}\).  The second follows by applying
  the decomposition of \ref{lem:chow-ii-algebraic-curvature-decomposition}.
\end{proof}

\begin{lemma}[quadratic with vanishing Weyl component]
  \label{lem:chow-ii-quadratic-zero-weyl}
  \source{chow-et-al-ricci-flow-part-ii:page-0127,chow-et-al-ricci-flow-part-ii:page-0128,chow-et-al-ricci-flow-part-ii:page-0129,chow-et-al-ricci-flow-part-ii:page-0130}
  \uses{lem:chow-ii-sharp-identity-formula, cor:chow-ii-trace-free-square-wedge-decomposition}
  \dcref{ch11:11.35}
  \group{Curvature Operator Algebra}
  Suppose \(R_W=0\), put
  \(\bar\lambda=\operatorname{Scal}(R)/n\) and
  \(\sigma=|\operatorname{Rc}_0(R)|^2/n\).  Then
  \[
  \begin{aligned}
  Q(R)={}&\frac1{n-2}\operatorname{Rc}_0\wedge\operatorname{Rc}_0
   +\frac{2\bar\lambda}{n-1}\operatorname{Rc}_0\wedge\operatorname{id}\\
   &-\frac2{(n-2)^2}(\operatorname{Rc}_0^2)_0\wedge\operatorname{id}
   +\left(\frac{\bar\lambda^2}{n-1}+\frac\sigma{n-2}\right)I.
  \end{aligned}
  \]
\end{lemma}
\begin{proof}
\sketch
  Write \(R=\bar\lambda I/(n-1)+R_0\), where
  \(R_0=2\operatorname{id}\wedge\operatorname{Rc}_0/(n-2)\).  By
  \ref{lem:chow-ii-sharp-identity-formula}, the scalar and mixed terms in
  \(Q(R)\) are the corresponding scalar and mixed terms displayed above.
  Diagonalize \(\operatorname{Rc}_0\), with eigenvalues \(\lambda_i\).  Then
  \(R_0(e_i\wedge e_j)=(\lambda_i+\lambda_j)e_i\wedge e_j/(n-2)\), while the
  intrinsic sharp formula gives
  \[
  R_0^\#(e_i\wedge e_j)=\sum_{k\ne i,j}
    \frac{(\lambda_i+\lambda_k)(\lambda_k+\lambda_j)}{(n-2)^2}
      e_i\wedge e_j.
  \]
  Using \(\sum_i\lambda_i=0\) and \(\sum_i\lambda_i^2=n\sigma\), the sum
  equals the eigenvalue of the remaining three terms in the formula.
\end{proof}

\begin{corollary}[Weyl and Ricci parts of the zero-Weyl quadratic]
  \label{cor:chow-ii-zero-weyl-quadratic-components}
  \source{chow-et-al-ricci-flow-part-ii:page-0127,chow-et-al-ricci-flow-part-ii:page-0128}
  \uses{lem:chow-ii-quadratic-zero-weyl, lem:chow-ii-ricci-contraction-wedge}
  \dcref{ch11:11.36}
  \group{Curvature Operator Algebra}
  If \(R_W=0\), then
  \[
  Q(R)_W=\frac1{n-2}(\operatorname{Rc}_0\wedge\operatorname{Rc}_0)_W
  \]
  and
  \[
  \operatorname{Rc}(Q(R))=-\frac2{n-2}(\operatorname{Rc}_0^2)_0
   +\frac{n-2}{n-1}\bar\lambda\operatorname{Rc}_0
   +(\bar\lambda^2+\sigma)\operatorname{id}.
  \]
\end{corollary}
\begin{proof}
\sketch
  Project the formula of \ref{lem:chow-ii-quadratic-zero-weyl} to the Weyl
  summand, then apply \ref{lem:chow-ii-ricci-contraction-wedge} and use
  \(\operatorname{Rc}_0^2=(\operatorname{Rc}_0^2)_0+\sigma\operatorname{id}\).
\end{proof}

\begin{theorem}[formula for the conjugation defect]
  \label{thm:chow-ii-conjugation-defect-formula}
  \source{chow-et-al-ricci-flow-part-ii:page-0119,chow-et-al-ricci-flow-part-ii:page-0120,chow-et-al-ricci-flow-part-ii:page-0121,chow-et-al-ricci-flow-part-ii:page-0122,chow-et-al-ricci-flow-part-ii:page-0123,chow-et-al-ricci-flow-part-ii:page-0124,chow-et-al-ricci-flow-part-ii:page-0125,chow-et-al-ricci-flow-part-ii:page-0126}
  \uses{def:chow-ii-conjugated-curvature-quadratic, def:chow-ii-algebraic-ricci-contraction}
  \dcref{ch11:11.27}
  \group{Equivariant Curvature Transformations}
  Let \(\operatorname{Rc}=\operatorname{Rc}(R)\),
  \(\bar\lambda=\operatorname{Scal}(R)/n\), and
  \(\sigma=|\operatorname{Rc}_0|^2/n\).  Whenever \(\ell_{a,b}\) is
  invertible,
  \[
  \begin{aligned}
  D_{a,b}(R)={}&((n-2)b^2-2(a-b))\operatorname{Rc}_0\wedge\operatorname{Rc}_0
  2a\operatorname{Rc}\wedge\operatorname{Rc}
  2b^2\operatorname{Rc}_0^2\wedge\operatorname{id}\\
  &+\frac{\sigma}{1+2(n-1)a}
  \bigl(nb^2(1-2b)-2(a-b)(1-2b+nb^2)\bigr)I.
  \end{aligned}
  \]
  In particular, \(D_{a,b}(R)\) depends only on \(\operatorname{Rc}(R)\).
  Moreover,
  \[
  \begin{aligned}
  \operatorname{Rc}(D_{a,b}(R))={}&-2b\operatorname{Rc}_0^2
  2(n-2)a\bar\lambda\operatorname{Rc}_0
  2(n-1)a\bar\lambda^2\operatorname{id}\\
  &+\frac{2(n-1)b+(n-2)^2b^2-2(n-1)a(1-2b)}{1+2(n-1)a}
     \sigma\operatorname{id}.
  \end{aligned}
  \]
\end{theorem}
\begin{proof}
\sketch
  \uses{lem:chow-ii-sharp-identity-formula,lem:chow-ii-quadratic-zero-weyl,cor:chow-ii-trace-free-square-wedge-decomposition}
  Polarizing \(D\), \ref{lem:chow-ii-sharp-identity-formula} shows that its
  mixed term with a Weyl operator vanishes; hence \(D(R+W)=D(R)\) for
  \(W\in\ker\operatorname{Rc}\).  We may therefore suppose \(R_W=0\).
  Apply \(Q\) to
  \[
  \ell_{a,b}(R)=\frac{1+2(n-1)a}{n-1}\bar\lambda I
    \frac{2(1+(n-2)b)}{n-2}\operatorname{id}\wedge\operatorname{Rc}_0,
  \]
  use \ref{lem:chow-ii-quadratic-zero-weyl}, and then apply
  \(\ell_{a,b}^{-1}\), which divides the scalar and trace-free Ricci pieces by
  the two displayed factors and fixes the Weyl piece.  Subtracting \(Q(R)\)
  and replacing the Weyl projection of
  \(\operatorname{Rc}_0\wedge\operatorname{Rc}_0\) by the identity in
  \ref{cor:chow-ii-trace-free-square-wedge-decomposition} gives the first
  formula after collecting coefficients.  Applying
  \ref{lem:chow-ii-ricci-contraction-wedge} to that formula gives the second.
\end{proof}

\begin{corollary}[eigenvalues of the conjugation defect]
  \label{cor:chow-ii-conjugation-defect-eigenvalues}
  \source{chow-et-al-ricci-flow-part-ii:page-0120,chow-et-al-ricci-flow-part-ii:page-0121}
  \uses{thm:chow-ii-conjugation-defect-formula}
  \dcref{ch11:11.29}
  \group{Equivariant Curvature Transformations}
  If \(\operatorname{Rc}_0e_i=\lambda_i e_i\), then \(e_i\wedge e_j\) is an
  eigenvector of \(D_{a,b}(R)\), with eigenvalue
  \[
  \begin{aligned}
  d_{ij}={}&((n-2)b^2-2(a-b))\lambda_i\lambda_j
   +2a(\lambda_i+\bar\lambda)(\lambda_j+\bar\lambda)
   +b^2(\lambda_i^2+\lambda_j^2)\\
   &+\frac{\sigma}{1+2(n-1)a}
   \bigl(nb^2(1-2b)-2(a-b)(1-2b+nb^2)\bigr).
  \end{aligned}
  \]
  The corresponding eigenvalue of \(\operatorname{Rc}(D_{a,b}(R))\) on
  \(e_i\) is obtained by applying its formula in
  \ref{thm:chow-ii-conjugation-defect-formula} to \(e_i\).
\end{corollary}
\begin{proof}
\sketch
  In the chosen basis, each of
  \(\operatorname{Rc}_0\wedge\operatorname{Rc}_0\),
  \(\operatorname{Rc}\wedge\operatorname{Rc}\),
  \(\operatorname{Rc}_0^2\wedge\operatorname{id}\), and \(I\) is diagonal on
  \(e_i\wedge e_j\).  Their eigenvalues are respectively
  \(\lambda_i\lambda_j\),
  \((\lambda_i+\bar\lambda)(\lambda_j+\bar\lambda)\),
  \((\lambda_i^2+\lambda_j^2)/2\), and \(1\).  Substitute these into the
  theorem; the Ricci assertion is immediate from its second formula.
\end{proof}
